
We characterize now when $\mcc_{\mathbb{G}}$ is a minimal resolution
of $\mbk[x]/I(\mcl)$. This result generalises the first main result in
\cite{ms}, namely Theorem 2 there. In a follow-up paper, we shall
investigate applications of our results, beginning with
generalisations of the remaining results of \cite{ms}, before
branching out more widely.

\begin{corollary}\label{min-resolution}
Let $\mbg$ be a strongly connected digraph or, equivalently, let $L$
be an $\icb$ matrix. Let $\mcc_{\mathbb{G}}$ be its associated $\cyc$
complex. Then
\begin{multline*}
\mcc_{\mathbb{G}}:\phantom{+} 0\leftarrow \mbk[x]/I(\mcl)\leftarrow
\mcc_0=\mbk[x]\xleftarrow{\sdif_1} \mcc_1=\mbk[x]^{r_1}
\xleftarrow{\sdif_2} \cdots \\ \cdots 
\xleftarrow{\sdif_{n-1}}\mcc_{n-1}=\mbk[x]^{r_{n-1}}\xleftarrow{} 0,
\end{multline*}
is a minimal free resolution of $\mbk[x]/I(\mcl)$ if and only if
$\mbg$ is strongly complete.
\end{corollary}
\begin{proof}
Suppose that $\mbg$ is strongly complete. That $\mcc_{\mathbb{G}}$ is
a minimal free resolution of $\mbk[x]/I(\mcl)$ follows from
Theorem~\ref{main} and Remark~\ref{dif-m}.

Conversely, suppose that $\mcc_{\mathbb{G}}$ is a minimal free
resolution of $\mbk[x]/I(\mcl)$. Consider $i,j\in [n-1]$ with $i\neq
j$, together with a basis element of the form 
\begin{eqnarray*}
(\{i\},\{j\},I_3,\ldots ,I_{k+1})\in\mcb_k.
\end{eqnarray*}
Then $\pm x^{a_{i,j}}_i$ appears as a coefficient in the expression
for $\dif_k(\{i\},\{j\},I_3,\ldots ,I_{k+1})$. By hypothesis,
$x^{a_{i,j}}_i\in\mfm$, and it follows that $a_{i,j}>0$. Similarly,
$a_{j,i}>0$.

Next consider a basis element of the form $(I_1,\ldots
,I_{k-1},\{i\},\{n\})\in\mcb_k$. A similar argument shows that
$a_{i,n}>0$.

Finally, consider a basis element of the form $(\{i\},I_2,\ldots
,I_k,\{n\})\in\mcb_k$. Again, a similar argument shows that
$a_{n,i}>0$, and the result follows.
\end{proof}
